\documentclass[11pt]{amsart}
\usepackage[T1]{fontenc}
\usepackage{lmodern}
\usepackage[margin=1.12in]{geometry}
\usepackage{amsmath,amssymb,amsthm}
\usepackage{microtype}
\usepackage[colorlinks=true,linkcolor=blue,citecolor=blue,urlcolor=blue]{hyperref}
\newtheorem{theorem}{Theorem}[section]
\newtheorem{lemma}[theorem]{Lemma}
\newtheorem{proposition}[theorem]{Proposition}
\newcommand{\R}{\mathbb R}
\newcommand{\N}{\mathbb N}
\newcommand{\Aut}{\operatorname{Aut}}
\newcommand{\ran}{\operatorname{ran}}
\newcommand{\ip}[2]{\langle #1,#2\rangle}
\title{Symmetric one-dependent colorings of regular trees}
\author{Kaiwen Jin}
\email{yui.ui.siki@gmail.com}
\date{October 13, 2026}
\subjclass[2020]{60G60, 05C15, 60G10}
\keywords{Finitely dependent process, proper coloring, regular tree, invariant law}
\begin{document}
\begin{abstract}
For every integer $d\geq3$, we construct a proper coloring of the infinite
$d$-regular tree with finitely many colors whose law is one-dependent and
invariant under every tree automorphism and every permutation of the colors.
The construction uses symmetric tensors in a real Hilbert space. Orthogonal
projections force adjacent colors to differ, tensor summation gives consistent
finite-dimensional laws, and a uniform estimate on rooted trees proves their
positivity. An explicit, unoptimized bound is
$q(d)=2^{2d+16}d^2(2^d-1)9^{d-1}$ colors.
\end{abstract}
\maketitle

\section{Introduction}

A random proper coloring has local exclusions: neighboring vertices must
receive different colors. Finite dependence asks for exact independence once
two sets of vertices are sufficiently separated, despite those exclusions.
Let $T_d$ be the infinite connected tree in which every vertex has degree $d$,
and let $[q]=\{1,\ldots,q\}$. A field $X=(X_z)_{z\in V(T_d)}$ taking values in
$[q]$ is \emph{one-dependent} if $X|_A$ and $X|_B$ are independent whenever
$A,B\subset V(T_d)$ are finite and their graph distance is greater than one.
Here independence refers to the complete vectors, including their internal
correlations. The field is proper if $X_z\ne X_w$ for every edge $zw$, almost
surely.

\begin{theorem}\label{thm:main}
For every integer $d\geq3$, put
\begin{equation}\label{eq:q}
q(d)=2^{2d+16}d^2(2^d-1)9^{d-1}.
\end{equation}
There is a proper one-dependent $[q(d)]$-valued coloring of $T_d$ whose law is
invariant under the full group $\Aut(T_d)$ and under every permutation of
$[q(d)]$.
\end{theorem}

The number of colors is fixed before any finite set is considered. The bound
in \eqref{eq:q} is chosen for a uniform positivity estimate and is not claimed
to be optimal.

Holroyd and Liggett constructed a stationary one-dependent four-coloring of
$\mathbb Z$ with a law invariant under color permutations
\cite[Theorems 1 and 4]{HL}. Liggett and Tang studied one-dependent colorings
of the three-ray star graph and asked about automorphism-invariant
one-dependent colorings of regular trees \cite{LT}. Their star has a single
branch point; a regular tree branches at every vertex. For unimodular
transitive graphs of bounded degree, Tim\'ar obtained automorphism-invariant two-dependent colorings
\cite[Theorem 2]{Timar}. Theorem~\ref{thm:main} gives one-dependence on regular
trees together with full color-permutation invariance.

Our construction is inspired by the word-space projections and positive-cone
argument for paths in Flin, Li and Suomela \cite[Section 4.1 and Appendix B]{FLS}.
We prove the estimates needed here directly. Branching requires a tensor with
one slot for each incident edge, rather than a succession of path operators.
Its sum over colors must allow deletion of an internal vertex, and positivity
must survive simultaneous contraction against several child branches. An
inverse sum of tensor projections supplies the deletion identity. A small
neighborhood of a product tensor supplies positivity on every finite tree.
The two parts of a tree bipartition use complementary projection ranges;
conjugation by word parity removes the apparent dependence on that choice.

Section~\ref{sec:projections} constructs the tensors and proves a quantitative
approximation by product tensors. Section~\ref{sec:finite} obtains positive,
consistent laws on finite forests. Section~\ref{sec:infinite} extends them to
the infinite tree and proves the required independence and symmetries.
All Hilbert spaces and contractions below are real.

\section{Projections on the word space}\label{sec:projections}

Fix $d\geq3$. Initially let $q$ be any integer satisfying
\begin{equation}\label{eq:invertibilityrange}
q\geq32d16^d.
\end{equation}
Let $G$ consist of the empty word and all finite words in $[q]$ whose adjacent
letters are distinct. On $H=\ell^2(G;\R)$ write $e_g$ for the coordinate unit
vector at a word $g$, and write $e$ for the unit vector at the empty word.
Define $L_i$ by prepending $i$ to a word unless its first letter is $i$, in
which case that letter is deleted. Define $R_i$ by the analogous operation at
the last letter. These operations permute the coordinate vectors.

Both operations are involutions. For a fixed $i$, the two end operations
commute. On words of length at least two this follows by performing the end
operations; on the empty word both composites return the empty word. On a
one-letter word $j$, both composites return $i$ if $j=i$, and $iji$ otherwise.
Thus $L_i$ and $R_i$ are commuting self-adjoint isometries. Set
\begin{equation}\label{eq:PQ}
P_i=\frac{(I+L_i)(I+R_i)}4,\qquad
J e_g=(-1)^{|g|}e_g,\qquad Q_i=JP_iJ.
\end{equation}
The operator $J$ is a self-adjoint isometry, $Je=e$, and $J$ anticommutes
with $L_i$ and $R_i$. Consequently
\[
Q_i=\frac{(I-L_i)(I-R_i)}4,\qquad P_iQ_i=0.
\]
Both $P_i$ and $Q_i$ are orthogonal projections. The two vectors
\begin{equation}\label{eq:uv}
u_i=e+e_i,\qquad v_i=e-e_i
\end{equation}
satisfy
\begin{equation}\label{eq:overlaps}
\|u_i\|=\|v_i\|=\sqrt2,\qquad
\ip e{u_i}=\ip e{v_i}=1,\qquad
\ip{u_i}{v_j}=\begin{cases}0&i=j,\\1&i\ne j.\end{cases}
\end{equation}
Here $e_i$ is the coordinate vector at the one-letter word $i$.

For $k\geq1$ we identify $H^{\otimes k}$ with $\ell^2(G^k;\R)$, and set
$H^{\otimes0}=\R$. Tensor products are therefore ordinary products of
coordinates. Contraction of a slot against $h\in H$ is a bounded map of norm
at most $\|h\|$, by Cauchy--Schwarz. These conventions also define every
partial contraction used below.

\begin{lemma}\label{lem:inverse}
The bounded operator
\[
A=\sum_{i=1}^q P_i^{\otimes d}
\]
on $H^{\otimes d}$ is invertible under \eqref{eq:invertibilityrange}, and
\begin{equation}\label{eq:inversebound}
\|A^{-1}\|\leq\frac{2\,4^d}{q}.
\end{equation}
\end{lemma}

\begin{proof}
Expand $P_i^{\otimes d}$ into its $4^d$ patterns, choosing in each coordinate
one of $I,L_i,R_i,L_iR_i$. Each pattern defines an involutive permutation
$\phi_i$ of $G^d$. Separate its fixed points from its off-diagonal entries.
After averaging the patterns, write $A=D+B$, where $D$ is diagonal and $B$
has zero diagonal. The all-identity pattern contributes $q/4^d$ to every
diagonal entry, whereas each summand of $A$ has diagonal entries at most one.
Thus
\begin{equation}\label{eq:D}
\frac{q}{4^d}\leq D(g,g)\leq q.
\end{equation}

For a fixed pattern, orient each nontrivial edge $g\leftrightarrow\phi_i(g)$
from smaller to larger total word length, using a fixed total order to break
ties. At any tuple, at most $2d$ labels $i$ can give a nontrivial neighbor of
no larger length. Indeed, such an $i$ must be the first or last letter of a
nonempty coordinate on which the pattern acts. Outside this set of labels,
a single end operation increases the length by one. The operation $L_iR_i$
increases the length of a nonempty word by two, and fixes the empty word.
If every active coordinate is empty and every active operation is $L_iR_i$,
the tuple is fixed, so it was already removed.

Let $M$ be the matrix of the oriented edges, retaining multiplicities if
different labels give the same edge. One of its row and column sums is at
most $q$, and the other is at most $2d$. For a finitely supported $f$, weighted
Cauchy--Schwarz gives
\[
\|Mf\|_2^2\leq
\left(\sup_x\sum_y M_{xy}\right)
\left(\sup_y\sum_x M_{xy}\right)\|f\|_2^2
\leq2dq\|f\|_2^2.
\]
It extends to $H^{\otimes d}$ by density. The off-diagonal matrix of this
pattern is $M+M^*$, so averaging all patterns proves
\begin{equation}\label{eq:B}
\|B\|\leq2\sqrt{2dq}.
\end{equation}
Equations \eqref{eq:D}--\eqref{eq:B} imply
$\|D^{-1}B\|\leq2\,4^d\sqrt{2d/q}\leq1/2$. The norm-convergent Neumann
series for $(I+D^{-1}B)^{-1}$ now gives \eqref{eq:inversebound} from
$A=D(I+D^{-1}B)$.
\end{proof}

Define
\begin{equation}\label{eq:tensors}
V=A^{-1}e^{\otimes d},\qquad t_{i,d}=P_i^{\otimes d}V.
\end{equation}
For $0\leq k\leq d$, obtain $t_{i,k}$ by contracting $d-k$ slots against $e$,
and set $s_{i,k}=J^{\otimes k}t_{i,k}$. The operators $A$ and $A^{-1}$ commute
with permutations of the $d$ slots. Hence $V$ and $t_{i,d}$ are symmetric
tensors, and the definition of $t_{i,k}$ does not depend on which slots are
retained. The tensors satisfy
\begin{align}
\sum_i t_{i,k}=\sum_i s_{i,k}&=e^{\otimes k},\label{eq:sum}\\
t_{i,k}\in\ran(P_i^{\otimes k}),\qquad
s_{i,k}&\in\ran(Q_i^{\otimes k}).\label{eq:range}
\end{align}
Contracting any slot of $t_{i,k}$ against $e$ gives $t_{i,k-1}$, and the same
holds for $s_{i,k}$ because $Je=e$. A permutation $\pi$ of $[q]$ acts on every
letter of $G$, fixes $e$, and conjugates $P_i$ to $P_{\pi(i)}$. It commutes
with $A$ after summation. Thus it carries $t_{i,k}$ and $s_{i,k}$ to the tensors
indexed by $\pi(i)$.

\begin{lemma}\label{lem:approximation}
Put $C_d=2\,4^d\sqrt{2^d-1}$ and $\delta=C_dq^{-3/2}$. For every $i$ and
$0\leq k\leq d$,
\begin{equation}\label{eq:approximation}
\|t_{i,k}-q^{-1}u_i^{\otimes k}\|\leq\delta,\qquad
\|s_{i,k}-q^{-1}v_i^{\otimes k}\|\leq\delta.
\end{equation}
\end{lemma}

\begin{proof}
Let $V_*=2^de^{\otimes d}/q$. Since $P_i e=u_i/2$,
\[
AV_*-e^{\otimes d}
=\frac1q\sum_{i=1}^q\sum_{\varnothing\ne S\subseteq[d]}
\bigotimes_{j=1}^d h_j(i,S),\qquad
h_j(i,S)=\begin{cases}e_i&j\in S,\\e&j\notin S.\end{cases}
\]
These $q(2^d-1)$ terms are distinct orthogonal coordinate vectors. The residual
therefore has norm $\sqrt{(2^d-1)/q}$. Lemma~\ref{lem:inverse} gives
$\|V-V_*\|\leq C_dq^{-3/2}$. Apply $P_i^{\otimes d}$ and contract the
discarded slots against $e$, operations of norm at most one. The corresponding
contraction of $P_i^{\otimes d}V_*$ is $q^{-1}u_i^{\otimes k}$. Applying the
isometry $J^{\otimes k}$ proves the second estimate.
\end{proof}

\section{Probabilities on finite forests}\label{sec:finite}

Choose either bipartition of $T_d$. For a finite vertex set $F$ and a specified
coloring $x\in[q]^F$, use its induced forest. At a vertex $z$ of degree $k$
in that forest, place $t_{x_z,k}$ if $z$ is in the first part and $s_{x_z,k}$
if it is in the second. Assign the tensor slots to incident edges in any
order and contract the slots at the two ends of each edge by the real inner
product. Let $p_F(x)$ be the product of the resulting scalars over components.
An isolated vertex uses a tensor with $k=0$, and $p_\varnothing=1$.

These finite-tree contractions are well-defined. Contracting a leaf vector
into its neighbor has norm at most the product of the two norms. Induction
bounds a full scalar contraction in absolute value by the product of all
vertex-tensor norms, and bounds an outgoing vector contraction in the same
way. The contraction is multilinear and continuous in the finitely many
tensors. Approximate them by finite-support tensors: their finite sums agree
in all leaf-contraction orders, and continuity gives order independence for
the original tensors. No infinite tree is contracted in this definition.

\begin{proposition}\label{prop:consistency}
The functions $p_F$ vanish on improper colorings and satisfy, for each
$z\in F$,
\begin{equation}\label{eq:deletion}
\sum_{a=1}^q p_F(x_{F\setminus\{z\}},x_z=a)
=p_{F\setminus\{z\}}(x_{F\setminus\{z\}}).
\end{equation}
They are invariant under color permutations and under tree automorphisms.
In particular, $\sum_{x\in[q]^F}p_F(x)=1$.
\end{proposition}

\begin{proof}
Cut a monochromatic edge with color $i$ at both ends and contract the branches
on its two sides. By \eqref{eq:range}, the remaining endpoint vectors lie in
the orthogonal ranges of $P_i$ and $Q_i$. Their inner product is zero.

To sum over a vertex color, \eqref{eq:sum} replaces its tensor by
$e^{\otimes k}$. Each of its incident edges then contracts $e$ into the tensor
at the other endpoint, producing the prescribed tensor with one fewer slot.
This proves \eqref{eq:deletion}, including when deletion of an internal vertex
splits a component into several branches. Repeated deletion proves the total
mass identity and consistency under every restriction of finite sets.

The action of a color permutation on the word space preserves inner products
and permutes the tensor family, so it preserves $p_F$. Switching the two
parts of the bipartition applies $J$ at both ends of every edge. Since $J$ is
an isometry, those two actions leave each inner product unchanged. Thus the
law is independent of the bipartition choice. Slot symmetry removes the
choice of incident-edge order. Every automorphism of the connected tree
preserves or swaps its two bipartition classes, proving the stated tree
invariance.
\end{proof}

It remains to prove that the scalars $p_F(x)$ are nonnegative. Orthogonality
and summation alone do not give this. We now take $q=q(d)$ from
\eqref{eq:q}; it satisfies \eqref{eq:invertibilityrange}. With the constant
$C_d$ of Lemma~\ref{lem:approximation}, this choice can be written as
\begin{equation}\label{eq:constants}
q=\bigl[64dC_d(3/2)^{d-1}\bigr]^2,\qquad
\eta=\frac1{16d},\qquad
\tau=q\delta(3/2)^{d-1}=\frac1{64d}=\frac\eta4.
\end{equation}

\begin{lemma}\label{lem:positivity}
For this value of $q$, every proper coloring of every finite induced forest
has $p_F(x)>0$.
\end{lemma}

\begin{proof}
First consider a finite rooted branch with one additional outgoing edge at
its root. Use one tensor slot for that edge, and contract all edges within
the branch, obtaining a message $m\in H$. We prove by induction that, for
a properly colored branch with root color $i$ in the first part,
\begin{equation}\label{eq:cone}
m=z\ell,\qquad z=\ip e m>0,\qquad
\ip e\ell=1,\qquad \|\ell-u_i\|\leq\eta.
\end{equation}
For a root in the second part replace $u_i$ by $v_i$. The message remains in
the corresponding projection range by \eqref{eq:range}.

Suppose the root has $r\leq d-1$ children. By induction write their messages
as $z_j\ell_j$ with $z_j>0$; the case $r=0$ uses empty products and requires
no induction hypothesis. A root in the first part has children in the second
part, with colors $i_j\ne i$. After dividing out these already positive
scalars, Lemma~\ref{lem:approximation} gives
\begin{equation}\label{eq:message}
\frac{m}{\prod_j z_j}=c u_i+\epsilon,\qquad
c=\frac1q\prod_{j=1}^r\ip{u_i}{\ell_j},\qquad
\|\epsilon\|\leq\delta(3/2)^r.
\end{equation}
Indeed, $\|\ell_j\|\leq\sqrt2+\eta<3/2$. From
$\ip{u_i}{v_{i_j}}=1$ and \eqref{eq:cone}, each factor in $qc$ is at least
$1-\sqrt2\eta$. Therefore
\begin{equation}\label{eq:mainterm}
qc\geq(1-\sqrt2\eta)^r\geq1-r\sqrt2\eta\geq\frac78.
\end{equation}
The scalar $b=\ip e{c u_i+\epsilon}$ is at least
$(7/8-\tau)/q>0$. Normalizing \eqref{eq:message} by $b$ yields
\[
\left\|\frac{c u_i+\epsilon}{b}-u_i\right\|
=\left\|\frac{\epsilon-\ip e\epsilon\,u_i}{b}\right\|
\leq\frac{(1+\sqrt2)\tau}{7/8-\tau}\leq4\tau=\eta.
\]
Thus \eqref{eq:cone} holds at the root, including the leaf case. At a root in
the second part, the same calculation uses $v_i$ and the overlaps
$\ip{v_i}{u_{i_j}}=1$.

Finally root a whole tree component, with no outgoing edge. The root can
have $r\leq d$ children. Its scalar contraction is at least
\begin{equation}\label{eq:rootpositive}
\frac{\prod_j z_j}{q}
\left[\frac78-q\delta(3/2)^d\right]>0.
\end{equation}
Here \eqref{eq:mainterm} still holds because $r\leq d$ and
$d\sqrt2\eta=\sqrt2/16<1/8$, while
$q\delta(3/2)^d=(3/2)\tau\leq3/128$. This also covers an isolated vertex.
The product over components is strictly positive. The same $q$ and cone
radius work for every finite size.
\end{proof}

\section{The infinite coloring}\label{sec:infinite}

\begin{proof}[Proof of Theorem~\ref{thm:main}]
Proposition~\ref{prop:consistency} and Lemma~\ref{lem:positivity} show that
$p_F$ is a probability mass function on $[q]^F$ for every finite vertex set
$F$, and that these functions are projective under restriction. The vertex
set of $T_d$ is countable, and the discrete state space $[q]$ is a Borel
space. The extension theorem \cite[Theorem 5.16, p.~92]{Kallenberg} therefore
gives a field $X$ having exactly these finite-dimensional laws.

If finite sets $A,B$ have distance greater than one, their induced union
has no edge between $A$ and $B$. Every component belongs entirely to one of
the two sets. The definition by component contractions gives, for all
$x\in[q]^A$ and $y\in[q]^B$,
\[
\mathbb P(X|_A=x,X|_B=y)=p_{A\cup B}(x,y)=p_A(x)p_B(y).
\]
This is independence of the complete vectors. Empty restrictions are
deterministic and cause no exception.

All finite-dimensional laws have the two invariances proved in
Proposition~\ref{prop:consistency}. Finite-dimensional distributions determine
the law on the product measurable space
\cite[Proposition 2.2, p.~25]{Kallenberg}. Applying this to each tree
automorphism and each color permutation gives invariance of the entire law.
The transformations are measurable since each transformed coordinate is a
finite-state function of one original coordinate. For every edge, the
probability of equal endpoint colors is zero. There are countably many
edges, so all edges are proper simultaneously almost surely.
\end{proof}

\subsection*{AI-use statement}
OpenAI Codex assisted with literature searches, development and checking of
the construction and proof, and preparation of this manuscript. Independent
mathematical review is requested.

\begin{thebibliography}{5}
\bibitem{HL}
A. E. Holroyd and T. M. Liggett,
\emph{Finitely dependent coloring},
Forum of Mathematics, Pi \textbf{4} (2016), e9.
\href{https://doi.org/10.1017/fmp.2016.7}{doi:10.1017/fmp.2016.7}.

\bibitem{LT}
T. M. Liggett and W. Tang,
\emph{One-dependent colorings of the star graph},
Annals of Applied Probability \textbf{33} (2023), no.~6A, 4341--4365.
\href{https://doi.org/10.1214/22-AAP1920}{doi:10.1214/22-AAP1920}.

\bibitem{Timar}
\'A. Tim\'ar,
\emph{Finitely dependent random colorings of bounded degree graphs},
arXiv:2402.17068v1 (2024).
\href{https://arxiv.org/abs/2402.17068v1}{arXiv:2402.17068v1}.

\bibitem{FLS}
M. Flin, L. Li and J. Suomela,
\emph{Quantum Advantage for Distributed Symmetry Breaking},
arXiv:2609.26788v1 (2026).
\href{https://arxiv.org/abs/2609.26788v1}{arXiv:2609.26788v1}.

\bibitem{Kallenberg}
O. Kallenberg,
\emph{Foundations of Modern Probability},
first edition, Springer, New York, 1997.
\href{https://doi.org/10.1007/b98838}{doi:10.1007/b98838}.
\end{thebibliography}
\end{document}
